% Diapositiva 31
%  El mono 
%=============================================
% Puntos 1, 2, 3: Por qué E(x,y,z,A0) no se puede visualizar completa,
% qué significa fijar z, y cortes bidimensionales.
%=============================================

\section{Visualización y Mapas de Calor}

\begin{frame}{Visualización de $E(x,y,z,A_0)$: El Problema Dimensional}

    \begin{columns}[T]

        \column{0.48\textwidth}

        \begin{block}{¿Por qué no se puede graficar directamente?}
            \scriptsize
            La función de error depende \textbf{simultáneamente} de cuatro parámetros:
            \[
                E = E(x,\, y,\, z,\, A_0)
            \]
            \begin{itemize}
                \item Una gráfica convencional solo admite \textbf{dos dimensiones}.
                \item Representar $E$ completa requeriría un espacio \textbf{5D}.
                \item Por tanto, la visualización directa \textbf{no es posible}.
            \end{itemize}
            \vspace{0.15cm}
            \textbf{Estrategia:} Se fijan $z$ y $A_0$ en sus valores óptimos y se analiza:
            \[
                E(x, y)\Big|_{z = z_{\text{real}},\; A_0 = A_0^*}
            \]
            Esto reduce el problema a \textbf{dos variables} visualizables.
        \end{block}

        \column{0.48\textwidth}

        \begin{block}{¿Qué representa cada corte?}
            \scriptsize
            \begin{itemize}
                \item Cada corte es una \textbf{fotografía} de la función de error a una profundidad fija.
                \vspace{0.15cm}
                \item Permite observar la \textbf{distribución espacial} del error en el plano $(x, y)$.
                \vspace{0.15cm}
                \item Variando $z$ se obtiene una \textbf{secuencia de cortes} que describe el comportamiento volumétrico de $E$.
            \end{itemize}
        \end{block}

        \vspace{0.2cm}

        \begin{block}{Idea Central}
            \centering
            \scriptsize
            Fijar $z$ y $A_0$ convierte un problema \textbf{5D} en uno \textbf{2D} visualizable.
        \end{block}

    \end{columns}

    \vspace{0.2cm}
    \begin{center}
        \footnotesize
        \textbf{Los cortes bidimensionales son la herramienta clave para explorar $E(x,y,z,A_0)$.}
    \end{center}

\end{frame}